\section{相关工作}
\label{sec:related}

\paragraph{手工设计的 kernel。}
现有 ML 框架，如 TensorFlow XLA~\cite{tensorflow_xla, Tensorflow}、PyTorch~\cite{pytorch} 与 TensorRT~\cite{tensorrt}，大多采用按算子逐核执行的方式，并依赖 GPU 专家为各个独立算子手工设计与实现 kernel。以 attention 为例，已经出现了大量针对特定架构特征或使用场景的专用实现，包括 FlashAttention~\cite{dao2023flash, tri2023flashdecoding, hong2024flashdecoding}、FasterTransformer~\cite{fastertransformer} 与 FlashInfer~\cite{ye2025flashinfer}。当前系统建立在一个高度碎片化的专用 kernel 生态上，这使得把整个推理流水线统一到单个、整体性的 mega-kernel 中变得非常困难。

\paragraph{ML 编译器。}
大量工作研究了如何通过 {\em 编译器} 为张量程序自动生成高性能 kernel。TVM~\cite{tvm, tvm_auto_tuner}、Ansor~\cite{ansor}、Triton~\cite{tillet2019triton} 以及其他系统~\cite{zheng2020flextensor, hagedorn2023graphene, feng2022tensorir, hu2024korch}，都建立在 Halide~\cite{halide, autohalide} 所提出的 algorithm-schedule 分离思想之上。另一条研究路线则采用 {\em superoptimization}，从高层规格中自动搜索高效实现~\cite{TASO, wang2021pet, zheng2023einnet, tensat, unger2022unity}。不过，这些编译器的核心目标仍是算子级优化，并不支持生成统一的 mega-kernel，也无法协调跨算子的整体执行。

\paragraph{Mega-kernel。}
已有的 mega-kernel 研究基本都依赖手工设计。例如，FlashDMoE 将 mixture-of-experts 计算和跨 GPU 通信融合进单个手写 mega-kernel~\cite{aimuyo2025flashdmoe}；Spector 等人则为 LLaMA-1B 手工开发了一个低延迟 mega-kernel~\cite{touvron2023llama, megakernel}。这类方案需要大量工程投入和深入的 GPU 专业知识，而且通常无法跨模型、跨 GPU 泛化。相比之下，\sys 提供了一条编译器驱动的自动化路径，能够把张量程序直接转换为高度优化的 mega-kernel，从而无需手工融合或手写 mega-kernel 实现。
